Within these self-contained subset reruns, per-electrode detection quality remained concentrated at the frontopolar and neighbouring anterior frontal electrodes rather than distributed uniformly across the frontal region (Figure~\ref{fig:exp1_subset_per_electrode}). On Internal, within the frontal subset's own rerun, Fp2, Fp1, AF4 and AF3 were the four strongest electrodes, reaching $F_1$ values of 85.28\%, 84.84\%, 77.04\% and 70.70\%, respectively. Performance then fell sharply at the two remaining frontal electrodes, F3 and F4, to $42.68\%$ and $41.00\%$, respectively. On Cao2018, the same four strongest electrodes, Fp1, Fp2, F3 and F4, made up the entire frontal subset, reaching $F_1$ values of 78.80\%, 77.05\%, 55.97\% and 54.64\%, respectively.
Although every one of these frontal electrodes scored lower in its own subset rerun than the same electrode did in the full-montage run (Section~\ref{sec:exp1_reference}), the ranking of electrodes by $F_1$ was unchanged from the full-montage run in both datasets, so restricting Stage~A/B screening to the frontal subset alone did not reorder the strongest frontal electrodes.

The central electrodes again showed a dataset-specific asymmetry within their own subset rerun. On Internal, C3 reached $F_1=58.59\%$ while C4 reached only $F_1=1.73\%$, so the pronounced C3--C4 asymmetry seen in the full-montage run persisted in the central subset's self-contained rerun. Wheras on Cao2018, C3 and C4 remained close together, at $27.67\%$ and $26.62\%$, respectively, so this asymmetry stayed specific to Internal. Parietal and occipital electrodes were much lower in their own subset reruns. The best-performing parietal electrode was CP1 at $1.92\%$ on Internal and CP4 at $15.63\%$ on Cao2018, and the best-performing occipital electrode was PO3 at $3.35\%$ on Internal and P3 at $10.04\%$ on Cao2018. Overall, the self-contained subset reruns preserved the same anterior-to-posterior spatial pattern observed under the full montage, with detection concentrated at the frontopolar and anterior frontal sites and much lower performance at the central, parietal and occipital electrodes.

\begin{figure}[ht]
    \centering
    \includegraphics[width=0.9\linewidth]{fig_exp1_subset_per_electrode.pdf}
    \caption{Per-electrode detection F1 within each region's own
        channel-subset rerun, coloured by scalp region. Each bar is one electrode
        scored inside that region's own self-contained subset run.}
    \label{fig:exp1_subset_per_electrode}
\end{figure}
